\documentclass[10pt,oneside]{report}
\usepackage{soc2style}
\begin{document}
\doctitle{Readiness and Gap Assessment}{Design, operation, evidence and governance scorecards}{SOC2-04}
\section*{Operating convention and boundaries}
An AICPA criterion identifies an outcome against which controls may be evaluated. The selected point-of-focus headings identify relevant characteristics; they are not discrete certification checklist items. Every owner assignment, technological mechanism and audit test in this suite is a proposed example to be tailored to the system description, commitments, risk profile, period, and service organization responsibilities.

\begin{itemize}
\item Always establish system boundaries, the in-scope trust services categories, service commitments and system requirements first.
\item For a security-only engagement, use the common criteria. For additional trust services categories, include the corresponding additional criteria, subject to properly justified applicability.
\item Document complementary user entity controls and subservice organizations; do not treat vendor assurance reports as replacements for management ownership.
\item Privacy points of focus marked [P] or [C] describe typical processor or controller relevance; assess actual responsibilities rather than treating the annotations as absolute.
\end{itemize}

\chapter{Purpose and engagement framing}
Readiness assessment is a consulting or management exercise, not an AICPA audit opinion. Establish whether the intended SOC 2 examination is Type 1 or Type 2, select the trust services categories, identify reporting period, define the system boundary, and align controls to commitments and system requirements. The AICPA source describes readiness services as consulting engagements that produce findings and recommendations rather than attestation opinions (p. 14).

\chapter{Assessment work plan}
\begin{enumerate}
\item Initiate: identify executive sponsor, system owner, target report type, scope and source systems.
\item Discover: catalog services, inventories, diagrams, data flows, vendor dependencies and documented policies.
\item Design review: for each relevant criterion, review risk, objective, implementation, owner, frequency, and dependencies.
\item Operating walkthrough: select actual executions from the intended observation period and validate reliable evidence.
\item Gap triage: distinguish criterion applicability, design gaps, operation gaps, evidence gaps and scope/ownership gaps.
\item Remediate and retest: assign work items, ensure management approval for risk acceptance, retest using independent records.
\item Prepare handoff: preserve dated evidence index, population definitions, unresolved issues, system description draft and practitioner questions.
\end{enumerate}

\subsection{Score dimensions -- a proposed management rubric}
\begin{itemize}
\item D: Design effectiveness: 0 = absent; 1 = informal; 2 = documented and reviewed; 3 = designed and risk-aligned with all key dependencies.
\item O: Operating effectiveness: 0 = not implemented; 1 = sporadic; 2 = executed with exceptions; 3 = sustained throughout evaluated period.
\item E: Evidence quality: 0 = none; 1 = assertions/screenshots only; 2 = attributable, scoped artifacts; 3 = complete and independently reproducible evidence.
\item G: Governance: 0 = no owner; 1 = informal ownership; 2 = accountable owner with procedure; 3 = reviewed accountability and timely remediation.
\end{itemize}

This 0--3 rubric is an internal readiness convention and has no AICPA endorsement. It is intentionally not used to claim that a numerical grade proves SOC 2 compliance. A criterion score must never substitute for a judgment about the actual control outcome.

\subsection{Worked readiness example}
Hypothetical CC8.1: protected main branch has peer review but CodeQL runs only after merge. Design might score D=1 or 2 depending on the promised pre-merge gate; operating observation O=2 for established PR review; evidence E=2 if change history is captured; governance G=2 if control ownership exists. Corrective action: set the scan as a required pre-merge check for applicable workflows, define emergency exception handling, and retest on a production-representative change. This example is not an audit conclusion.

\chapter{Complete criterion assessment register}
Use the full 61-criterion register to record design, execution, evidence and governance scores. Mark out-of-scope or not-applicable determinations separately with a reason; never silently omit a criterion from a selected category. The rows below are a working worksheet and not a representation of AICPA prescribed testing steps.

\small\begin{longtable}{@{}p{0.63in}p{2.13in}p{2.28in}p{0.92in}@{}}
\toprule ID & Control domain / owner & Primary starting evidence & D/O/E/G \\ \midrule\endfirsthead
\toprule ID & Control domain / owner & Primary starting evidence & D/O/E/G \\ \midrule\endhead
\textbf{CC1.1} & Ethics and conduct / Legal / HR & Approved conduct policy; attestations; case register with redactions & \_ / \_ / \_ / \_ \\
\addlinespace[3pt]
\textbf{CC1.2} & Board oversight / Board secretary & Board or committee minutes; risk packs; action tracker & \_ / \_ / \_ / \_ \\
\addlinespace[3pt]
\textbf{CC1.3} & Structures and responsibility / GRC / Executive & Org chart; role descriptions; RACI; delegations & \_ / \_ / \_ / \_ \\
\addlinespace[3pt]
\textbf{CC1.4} & Competence and training / HR / Security & Role matrix; training records; screening attestations & \_ / \_ / \_ / \_ \\
\addlinespace[3pt]
\textbf{CC1.5} & Accountability and performance / HR / Compliance & Review rubric; acknowledgments; corrective-action logs & \_ / \_ / \_ / \_ \\
\addlinespace[3pt]
\textbf{CC2.1} & Quality information and inventory / Asset / Data Governance & CMDB; data-flow diagram; classification register; reconciliation logs & \_ / \_ / \_ / \_ \\
\addlinespace[3pt]
\textbf{CC2.2} & Internal communications / GRC / Comms & Communications log; trainings; incident channel instructions & \_ / \_ / \_ / \_ \\
\addlinespace[3pt]
\textbf{CC2.3} & External communications / Legal / Customer Trust & Service agreements; trust communications; contact playbook & \_ / \_ / \_ / \_ \\
\addlinespace[3pt]
\textbf{CC3.1} & Objectives and risk appetite / Executive / Risk & Risk appetite; system description; SLA register & \_ / \_ / \_ / \_ \\
\addlinespace[3pt]
\textbf{CC3.2} & Risk identification and analysis / Enterprise Risk & Risk register; workshop results; treatment approvals & \_ / \_ / \_ / \_ \\
\addlinespace[3pt]
\textbf{CC3.3} & Fraud risk / Risk / Internal Audit & Fraud scenarios; SoD review; whistleblower handling & \_ / \_ / \_ / \_ \\
\addlinespace[3pt]
\textbf{CC3.4} & Significant change assessment / Architecture / Risk & Change intake; impact assessments; risk decisions & \_ / \_ / \_ / \_ \\
\addlinespace[3pt]
\textbf{CC4.1} & Control monitoring / Control Assurance & Control monitoring schedule; dashboards; findings & \_ / \_ / \_ / \_ \\
\addlinespace[3pt]
\textbf{CC4.2} & Deficiency communication / GRC / Risk & Findings log; escalation tickets; validation artifacts & \_ / \_ / \_ / \_ \\
\addlinespace[3pt]
\textbf{CC5.1} & Risk-based control activities / Risk / Control Owners & Risk-control matrix; approvals; control narratives & \_ / \_ / \_ / \_ \\
\addlinespace[3pt]
\textbf{CC5.2} & Technology general controls / IT / Platform & ITGC inventory; platform standards; administrative review & \_ / \_ / \_ / \_ \\
\addlinespace[3pt]
\textbf{CC5.3} & Policy deployment / GRC / Operational Owners & Policies; procedures; schedules; exception register & \_ / \_ / \_ / \_ \\
\addlinespace[3pt]
\textbf{CC6.1} & Logical access architecture / IAM / Platform & IAM configuration; architecture; identity inventory & \_ / \_ / \_ / \_ \\
\addlinespace[3pt]
\textbf{CC6.2} & Provisioning and deprovisioning / IAM / HR & Access requests; approval log; HR termination feed & \_ / \_ / \_ / \_ \\
\addlinespace[3pt]
\textbf{CC6.3} & Least privilege review / IAM / Application Owners & Access review signoffs; role policy; permissions exports & \_ / \_ / \_ / \_ \\
\addlinespace[3pt]
\textbf{CC6.4} & Physical access / Facilities / Cloud Provider & Badge access review; facility records; subservice reports & \_ / \_ / \_ / \_ \\
\addlinespace[3pt]
\textbf{CC6.5} & Secure asset disposal / IT Asset Management & Sanitization records; disposal certificates; asset ledger & \_ / \_ / \_ / \_ \\
\addlinespace[3pt]
\textbf{CC6.6} & Perimeter and external threat defenses / Network / Cloud & Firewall rules; segmentation map; WAF/config reviews & \_ / \_ / \_ / \_ \\
\addlinespace[3pt]
\textbf{CC6.7} & Secure data movement / Data Security & DLP policy; TLS configurations; export approvals & \_ / \_ / \_ / \_ \\
\addlinespace[3pt]
\textbf{CC6.8} & Malicious software prevention / Endpoint / Platform & EDR coverage; quarantine events; allowlist policy & \_ / \_ / \_ / \_ \\
\addlinespace[3pt]
\textbf{CC7.1} & Vulnerability and drift detection / Vulnerability Management & Scan results; configuration diffs; risk acceptance tickets & \_ / \_ / \_ / \_ \\
\addlinespace[3pt]
\textbf{CC7.2} & Anomaly monitoring / SOC / SRE & Detection rules; SIEM alerts; on-call histories & \_ / \_ / \_ / \_ \\
\addlinespace[3pt]
\textbf{CC7.3} & Security event evaluation / SOC / Incident Response & Triage cases; severity rubric; incident timeline & \_ / \_ / \_ / \_ \\
\addlinespace[3pt]
\textbf{CC7.4} & Incident response / Security Operations & IR plan; tabletop results; incident cases & \_ / \_ / \_ / \_ \\
\addlinespace[3pt]
\textbf{CC7.5} & Security incident recovery / SRE / Business Continuity & Recovery playbooks; restore tests; PIR action list & \_ / \_ / \_ / \_ \\
\addlinespace[3pt]
\textbf{CC8.1} & Controlled change management / Engineering / CAB & Pull requests; checks; releases; change tickets & \_ / \_ / \_ / \_ \\
\addlinespace[3pt]
\textbf{CC9.1} & Business disruption mitigation / BCM / Risk & BIA; continuity response plan; mitigation decisions & \_ / \_ / \_ / \_ \\
\addlinespace[3pt]
\textbf{CC9.2} & Third-party risk / Third-Party Risk & Vendor register; diligence; contracts; monitoring evidence & \_ / \_ / \_ / \_ \\
\addlinespace[3pt]
\textbf{A1.1} & Capacity management / SRE / Capacity & Capacity dashboards; forecast; scaling decisions & \_ / \_ / \_ / \_ \\
\addlinespace[3pt]
\textbf{A1.2} & Environmental and backup protections / Infrastructure / BCM & Backup jobs; retention settings; facility/provider evidence & \_ / \_ / \_ / \_ \\
\addlinespace[3pt]
\textbf{A1.3} & Recovery testing / BCM / SRE & DR exercise plan; restore logs; issue remediation & \_ / \_ / \_ / \_ \\
\addlinespace[3pt]
\textbf{C1.1} & Confidential information life cycle / Data Governance & Classification inventory; data access matrix; retention map & \_ / \_ / \_ / \_ \\
\addlinespace[3pt]
\textbf{C1.2} & Confidential disposal / Data Governance / IT & Retention schedule; deletion jobs; vendor certificates & \_ / \_ / \_ / \_ \\
\addlinespace[3pt]
\textbf{PI1.1} & Processing specifications / Product / Data Engineering & Data dictionary; schema registry; acceptance criteria & \_ / \_ / \_ / \_ \\
\addlinespace[3pt]
\textbf{PI1.2} & Input completeness and accuracy / Data Engineering & Validation rules; rejects; ingestion logs & \_ / \_ / \_ / \_ \\
\addlinespace[3pt]
\textbf{PI1.3} & Processing correctness / Platform / Engineering & Job runs; test results; reconciliation outcomes & \_ / \_ / \_ / \_ \\
\addlinespace[3pt]
\textbf{PI1.4} & Output distribution / Product / Data Engineering & Delivery manifests; approvals; integrity checks & \_ / \_ / \_ / \_ \\
\addlinespace[3pt]
\textbf{PI1.5} & Processing records retention / Data Engineering / Records & Retention policies; transaction logs; storage tests & \_ / \_ / \_ / \_ \\
\addlinespace[3pt]
\textbf{P1.1} & Privacy notice / Privacy / Legal & Privacy notice versions; publication history; data map & \_ / \_ / \_ / \_ \\
\addlinespace[3pt]
\textbf{P2.1} & Choice and consent / Privacy / Product & Consent ledger; UX screens; withdrawal logs & \_ / \_ / \_ / \_ \\
\addlinespace[3pt]
\textbf{P3.1} & Purpose-limited collection / Privacy / Product & Data inventory; form design; API schemas & \_ / \_ / \_ / \_ \\
\addlinespace[3pt]
\textbf{P3.2} & Explicit consent collection / Privacy / Legal & Consent artifacts; notices; audit logs & \_ / \_ / \_ / \_ \\
\addlinespace[3pt]
\textbf{P4.1} & Purpose limitation / Privacy / Data Owners & Purpose map; access policy; data-use approvals & \_ / \_ / \_ / \_ \\
\addlinespace[3pt]
\textbf{P4.2} & Personal-information retention / Privacy / Records & Retention schedule; legal holds; purge reports & \_ / \_ / \_ / \_ \\
\addlinespace[3pt]
\textbf{P4.3} & Personal-information disposal / Privacy / Data Engineering & Erasure tickets; deletion evidence; vendor confirmations & \_ / \_ / \_ / \_ \\
\addlinespace[3pt]
\textbf{P5.1} & Data subject access / Privacy Operations & DSAR register; identity check; response record & \_ / \_ / \_ / \_ \\
\addlinespace[3pt]
\textbf{P5.2} & Correction and amendment / Privacy Operations / Data & Correction tickets; updated records; vendor notices & \_ / \_ / \_ / \_ \\
\addlinespace[3pt]
\textbf{P6.1} & Authorized disclosure / Privacy / Legal & Data-sharing inventory; agreements; consent reference & \_ / \_ / \_ / \_ \\
\addlinespace[3pt]
\textbf{P6.2} & Disclosure accounting / Privacy / Data Engineering & Disclosure log; export records; reconciliation & \_ / \_ / \_ / \_ \\
\addlinespace[3pt]
\textbf{P6.3} & Unauthorized disclosure recording / Privacy / Incident Response & Incident register; disclosure timeline; investigation records & \_ / \_ / \_ / \_ \\
\addlinespace[3pt]
\textbf{P6.4} & Privacy vendor commitments / Privacy / Third-Party Risk & DPA/contract clauses; vendor assessments; remediation log & \_ / \_ / \_ / \_ \\
\addlinespace[3pt]
\textbf{P6.5} & Vendor incident notifications / Privacy / Third-Party Risk & Notification terms; escalation contacts; vendor tests & \_ / \_ / \_ / \_ \\
\addlinespace[3pt]
\textbf{P6.6} & Privacy incident notification / Privacy / Legal & Breach decision tree; counsel review; notices & \_ / \_ / \_ / \_ \\
\addlinespace[3pt]
\textbf{P6.7} & Accounting and disclosure to subject / Privacy Operations & DSAR package; inventory; disclosure accounting & \_ / \_ / \_ / \_ \\
\addlinespace[3pt]
\textbf{P7.1} & Personal data quality / Data Governance & Data-quality dashboard; reconciliations; corrected records & \_ / \_ / \_ / \_ \\
\addlinespace[3pt]
\textbf{P8.1} & Privacy inquiry and dispute handling / Privacy / Compliance & Case register; SLA results; corrective actions & \_ / \_ / \_ / \_ \\
\addlinespace[3pt]
\bottomrule\end{longtable}\normalsize
\chapter{Gap findings and remediation register}
\subsection{Standardized issue template}
\begin{itemize}
\item Finding ID and classification: design / operating / evidence / ownership / scoped applicability.
\item Criterion reference(s), source requirement excerpt or reference, affected system/service and risk statement.
\item Population and period tested, sample identifiers, observed condition, expected control outcome and root-cause hypothesis.
\item Impact on service commitments and other relevant categories; management and practitioner communications.
\item Corrective action, owner, planned completion, interim mitigation, risk acceptance approval, retest outcome and closure.
\end{itemize}

\subsection{Prioritization model -- illustrative}
Prioritize gaps based on customer impact, exploitability, likelihood of recurrence, sensitivity of affected data, public exposure, incident indicators and proximity to the expected reporting period. High-severity exceptions need explicit escalation; a scoring model cannot override applicable legal obligations or a practitioner's engagement decisions.

\subsection{Readiness exit conditions}
\begin{enumerate}
\item Scope and system description reconciled to production reality and provider dependencies.
\item All selected-category criteria mapped to documented controls or rationalized applicability decisions.
\item Control owners have demonstrated at least one representative operation; Type 2 preparation additionally supports sustained operation over the intended reporting period.
\item Evidence sources support trustworthy populations, completeness, accuracy and preserved timestamps.
\item Known exceptions evaluated and assigned to accountable resolution with documented residual risk.
\item Practitioner-facing evidence index and management representations reviewed for consistency.
\end{enumerate}


\chapter*{Source and editorial protocol}
\addcontentsline{toc}{chapter}{Source and editorial protocol}
This independent document is grounded in the user-supplied \emph{AICPA TSP Section 100: 2017 Trust Services Criteria for Security, Availability, Processing Integrity, Confidentiality, and Privacy (With Revised Points of Focus---2022)}, 75-page PDF, especially paragraphs .01--.29, the criterion table on pp.~14--67, and the glossary. Criterion wording is quoted when labeled as such. Selected point-of-focus labels are aids to interpretation and may be incomplete. Engineering recommendations, role assignments, tools, tests, dates, scorecards, and example frequencies are author-developed examples and are not AICPA requirements. The 2022 revisions changed points of focus, not the 2017 criteria (source p.~3). The source states that an assessment of each point of focus is not required (p.~5). This document does not constitute a CPA attestation or legal advice.
\end{document}
